\documentclass{article}
\usepackage[left=2cm, right=2cm, top=2cm, bottom=2cm]{geometry}
\usepackage{graphicx}
\usepackage{amsmath}
\usepackage{amssymb}
\usepackage{amsthm}
\usepackage{fancyhdr}
\usepackage{verbatim}
\usepackage{listings}
\usepackage{xcolor}
\usepackage{pdfpages}
\usepackage{cancel}
\usepackage{physics}
\usepackage{esint}
\usepackage{braket}

\lstset{
    language=C++,
    basicstyle=\ttfamily\footnotesize,
    keywordstyle=\color{blue},
    commentstyle=\color{green},
    stringstyle=\color{red},
    numbers=left,
    numberstyle=\tiny\color{gray},
    frame=single,
    breaklines=true,
    captionpos=b,
}


\title{PHYS 427: Lecture \# 8}
\author{Cliff Sun}

\newtheorem{theorem}{Theorem}[section]
\newtheorem{lemma}[theorem]{Lemma}
\newtheorem{definition}[theorem]{Definition}
\newtheorem{conjecture}[theorem]{Conjecture}
\newtheorem{proposition}[theorem]{Proposition}
\newtheorem{corollary}[theorem]{Corollary}
\newtheorem{one minute paper}[theorem]{One Minute Paper}

\pagestyle{fancy}
\lhead{\textbf{\thepage}\ \ \nouppercase{\rightmark}}
\chead{PHYS 427: Lecture \# 8}
\rhead{Cliff Sun}

\begin{document}

\maketitle

\section*{Lecture Span}
\begin{enumerate}
    \item Heat and Work
\end{enumerate}

\section*{Recap}

Micro canonical ensemble: fixed energy, number of particles. Canonical ensemble has a reservoir, combined energy and temperature are constant. 

For MCE, $S(U,V)$ is maximized in equilibrium. For CE, $F(U,V)$ is minimized in equilibrium. 

For MCE, $\Omega(U) = \sum_{\text{all states with energy $U$}} 1$. 

For CE, $Z(T) = \sum e^{-\epsilon_i/k_B T}$. 

For MCE, $S = k_B \ln \Omega$. 

For CE, $F = -k_B T \ln Z$. 

For MCE, 

\begin{align}
    dS &= \left( \partial_{U}S \right)_V dU + \left( \partial_{V}S \right)dV \\
    &= \frac{1}{T}dU + \frac{p}{T}dU
\end{align}

For CE, 

\begin{align}
    dF &= \left( \partial_{T}F \right)_TdT  + \left( \partial_V F\right)_T dV \\
    &= -S dT - pdV
\end{align}

You can match definitions to then rewrite $T$ in terms of $S$, $p$ in terms of $F$, and etc. 

\section*{Heat and Work}

Thermodynamic identity:

\begin{equation}
    dU = TdS - pdV
\end{equation}

\begin{enumerate}
    \item Work: can be electrical, or mechanical
    \item Heat: cannot be completely converted into work because of the 2nd law of thermodynamics
\end{enumerate}

The 2nd law of thermodynamics states $\Delta S \geq 0$, and this is the \underline{total entropy of a system!}

Consider a reservoir of temperature $T_H$, then it takes heat $Q_H$, then outputs $Q_C$ (heat that goes into the cold reservoir) and $W$. Then 

\begin{equation}
    Q_H = W + Q_C
\end{equation}

Then 

\begin{align}
    \Delta S_{tot} &= \Delta S_{hot} + \Delta S_{cold} \geq 0 \\
    &= -\frac{Q_H}{T_H} + \frac{Q_C}{T_C} \geq 0
\end{align}

The bottom reservoir then must be colder than the top reservoir. Define efficiency, 

\begin{equation}
    \eta = \frac{W}{Q_H} = 1 - \frac{Q_C}{Q_H}
\end{equation}

The maximum possible efficiency for $\Delta S_{tot} = 0$, then 

\begin{align}
    \frac{Q_C}{Q_H} &= \frac{T_C}{T_H}
\end{align}

Then, the max efficiency is 

\begin{equation}
    \eta_{\max} = 1 - \frac{T_C}{T_H}
\end{equation}

This is the Carnot efficiency. But you can have losses along the way, such as $\Delta S_{tot} > 0$, which can be caused by a heat leak from H to C. 

\subsection*{Irreversible processes}

Create $\Delta S_{tot}>0$ that cannot be removed due to the 2nd law of thermodynamics. 

\section*{Carnot Cycle}

Idealized way to realize a heat engine at the maximal efficiency $\eta_{\max}$. Start at $(T_H, S_1)$, then heat up to $(T_H, S_2)$. Then go $(T_H, S_2)$ to $(T_C, S_2)$. Then to $(T_C, S_1)$, then close the loop. 

Here, 

\begin{equation}
    \oint dQ = (T_H - T_C)(S_2 - S_1)
\end{equation}

The area of this rectangle is the work done by the engine. You can then calculate 

\begin{equation}
    \eta = \frac{W}{Q_H} = \frac{(T_H - T_C)(S_2 - S_1)}{T_H(S_2 - S_1)} = 1 - \frac{T_C}{T_H} 
\end{equation}

Note, reversible cooling and heating ($\Delta S =0$) need to be very slow; otherwise, you gain entropy. But this is not practical. 

\subsection*{Monatomic ideal gas as working substance}

Sackur-Tetrode equation 

\begin{equation}
    S = Nk_B T\left( \ln\left( \frac{n_Q V}{N} \right) + \frac{5}{2} \right)
\end{equation}

$PV = Nk_B T$. for isothermal, $PV = \text{const.}$. 

For isentropcic, 

\begin{equation}
    S = Nk_B \left[ \ln(T^{3/2}) + \ln(V) \right]
\end{equation}

Then $T^{3/2}V$ is constant, which means that $p^{3/2}V^{5/2} = \text{const.}$. Then 

\begin{equation}
    pV^{\gamma} = \text{const.}
\end{equation}

For monatomic ideal gas. For $\gamma = 5/3 = C_P/C_V$. For an ideal gas, you can plot it on a PV diagram. 

First, isothermal expansion. Then, isentropic expansion. Then isothermal. Then isentropic. You can then calculate the work done by the system.

\begin{equation}
    W = (T_H - T_C)(S_2 - S_1)
\end{equation}

Irreversible expansion. Box with wall with door in the middle. Volume on left is $V_i$ and on right is $V_f$. $Q = 0$. Also, $\Delta U = 0 \iff \Delta T =0$.
The entropy $\Delta S = Nk_B\left[ \ln\left( V_f/V_i \right) \right] > 0$. For irreversible processes, $T dS > dQ$. For reversible, they are equal. 

\end{document}